\title{Large Language Models Can Automatically Engineer Features for Few-Shot Tabular Learning}

\begin{abstract}
Large Language Models (LLMs) have demonstrated exceptional proficiency in solving intricate and novel reasoning challenges, indicating substantial promise for tabular learning, which is essential to numerous practical applications. In this study, we introduce an innovative in-context learning paradigm named \textsf{FeatLLM}. This approach leverages LLMs as feature generators, crafting transformed data sets that are better optimized for tabular prediction tasks. The engineered features allow a simple downstream machine learning model—such as linear regression—to estimate class probabilities efficiently, delivering high accuracy in few-shot learning scenarios. Notably, the \textsf{FeatLLM} methodology relies solely on these generated features and a basic predictive model during inference, bypassing the need for recurrent LLM queries for each new data point. Furthermore, \textsf{FeatLLM} requires nothing more than API-level LLM access, thus circumventing issues with prompt length constraints. Experimental results across diverse tabular datasets from multiple domains confirm that \textsf{FeatLLM} consistently yields high-quality predictive rules and significantly (by an average margin of 10\%) surpasses methods such as TabLLM and STUNT.
\end{abstract}

\section{Introduction}
Large Language Models (LLMs) have achieved remarkable progress recently, exhibiting strong capabilities to generate contextually relevant outputs for unseen problems even without being specifically fine-tuned for the task~\cite{creswell2022selection,imani2023mathprompter,taylor2022galactica}. Owing to their training on vast text datasets, LLMs possess extensive prior knowledge that can effectively substitute for the expertise of human specialists in a wide range of disciplines~\cite{Genomes2021,singhal2023large,wilcox2003role}, thereby advancing automation in numerous fields such as dialogue analysis~\cite{finch2023leveraging} and program repair~\cite{fan2023automated}. Notably, when prompts include a concise task description and a few examples, LLMs are capable of discerning relevant context and emphasizing critical features and samples~\cite{brown2020language,zhao2023survey}. Such context-aware data handling highlights the potential for LLMs to assume the role of advanced feature engineers.

Recent efforts have extended LLMs’ capacity for knowledge transfer and reasoning to the area of tabular learning. For example, encoding prior information—such as an increased likelihood of diseases in elderly populations—has proven advantageous for predictive tasks in healthcare. In these applications, researchers commonly encode both feature and task descriptions in natural language, transform the tabular data accordingly, and provide these inputs to an LLM~\cite{dinh2022lift,wang2023anypredict}. This workflow is typically supplemented with in-context learning~\cite{nam2023semi} or parameter-efficient adaptation methods~\cite{dinh2022lift,hegselmann2023tabllm}. Empirical studies show that utilizing the inherent knowledge of LLMs yields performance improvements over standard tabular learning approaches, particularly in scenarios with limited labeled data~\cite{hegselmann2023tabllm}.

Prevailing LLM-based strategies for tabular learning are subject to several drawbacks. Primarily, in end-to-end prediction scenarios, each data instance necessitates at least one LLM inference, leading to considerable computational demand. Furthermore, achieving superior accuracy typically involves fine-tuning the LLM, which proves unfeasible for models that lack open-source training code or tuning APIs—a limitation especially salient given that many leading-edge LLMs now only offer restricted access via inference APIs~\cite{anil2023palm,openai2023gpt4}. Another major drawback is that numerous current methods perform poorly or become impractical when handling lengthy prompts~\cite{liu2023lost}; as the number of features within tabular data increases, thus lengthening serialized prompts, performance may degrade or inference may be rendered unmanageable. Collectively, these issues create significant hurdles for deploying LLM-based tabular learning systems in practice.

Our methodology departs from using LLMs solely for direct predictive tasks. Rather, we position LLMs within the feature engineering phase, thus capitalizing more effectively on their encoded prior knowledge. In particular, we introduce a distinct in-context learning approach, \textsf{FeatLLM}, focused on elucidating the “criteria” underpinning LLM-driven predictions. For example, when tasked with predicting the presence of a specific disease, the LLM can generate interpretable rules specifying the combinations of features that suggest the disease is present. These generated criteria or rules serve as new features, effectively replacing original attributes for subsequent learning tasks. By redirecting the LLM’s utility toward feature construction, we enable downstream application of simpler models—diminishing reliance on intricate inference schemes and thereby accelerating inference speeds. The power of in-context learning in LLMs permits the extraction of valuable features merely by incorporating illustrative demonstrations within prompts, thus bypassing the necessity for additional model training. Moreover, by drawing upon ensemble methodologies such as feature bagging~\cite{breiman2001random}, our approach can better manage the complications introduced by extensive prompt sizes.

\textsf{FeatLLM} decomposes the process of feature engineering within input prompts into two main sub-tasks: (1) grasping the underlying problem and discerning how features relate to targets; and (2) leveraging both the insights from this understanding and a set of few-shot examples to formulate discriminative rules that can separate different classes. This carefully segmented reasoning mechanism enables LLMs to disregard features that lack informativeness during rule extraction. The inferred rules are implemented directly on the dataset (interpreted as executable code), converting the data into binary indicators reflecting compliance with each respective rule. Subsequently, a simple model is trained to predict class probabilities using these synthesized features. To enhance resilience, an ensemble methodology is adopted, involving repeated feature generation and feature bagging. By strategically selecting the number of features and data samples for bagging, the issue of oversized prompts is alleviated. Operating without reliance on additional model training and offering efficient inference, \textsf{FeatLLM} permits the practical deployment of LLMs in a variety of application domains.

We assess \textsf{FeatLLM} across 13 tabular datasets under limited data conditions, highlighting its consistent and strong results. The experiments demonstrate that LLMs can effectively integrate existing domain knowledge with new insights drawn directly from the data, leading to the extraction of valuable features. Our framework consistently outperforms state-of-the-art few-shot learning baselines in multiple scenarios. The implementation is publicly accessible at the anonymized GitHub repository:~\texttt{https://github.com/Sungwon-Han/FeatLLM}.

\section{Related Work}

\subsection{Few-Shot Learning with Tabular Data} The progress in few-shot learning techniques has facilitated the derivation of robust, general-purpose representations from datasets, even when only limited annotation is available~\cite{chen2018closer}. While few-shot methods were initially tailored for image-based applications, their capacity for generalization has since been demonstrated for a range of data modalities, such as text, audio, and, more recently, tabular formats~\cite{majumder2022few,nam2023stunt,schick2021exploiting}. Nevertheless, relying on sparse labeled samples heightens susceptibility to capturing misleading or non-generalizable patterns~\cite{bartlett2021deep}. To mitigate this, contemporary research has leveraged supplementary information or integrated domain-relevant prior knowledge to enforce helpful inductive biases during training. Strategies include employing substantial quantities of unlabeled data—paired with suitable augmentation schemes—for semi-supervised learning~\cite{ucar2021subtab,yoon2020vime}, or seeking robust initial model parameters via unsupervised, task-agnostic pre-training~\cite{somepalli2021saint}. These unsupervised approaches may employ reconstruction objectives by masking or corrupting portions of data~\cite{arik2021tabnet,majmundar2022met}, or introduce contrastive learning via augmented data samples~\cite{bahri2022scarf,chen2020simple}. However, the intrinsic heterogeneity of tabular datasets complicates the design of universally effective augmentation procedures, thereby limiting the efficacy of such strategies—particularly in few-shot contexts~\cite{nam2023stunt}.

More recent investigations suggest pre-training on a wide assortment of tabular datasets, far beyond those directly related to a target problem, followed by transfer learning for downstream applications~\cite{zhang2023generative,levin2022transfer,zhu2023xtab}. For example, the TransTab framework~\cite{wang2022transtab} utilizes transformer architectures to capture semantic relations among columns—leveraging column names rather than raw cell values. Accumulated knowledge across various tables and columns allows models to enable zero- and few-shot prediction abilities on new tabular tasks. Separately, TabPFN~\cite{hollmann2022tabpfn} synthesizes pre-training data by blending diverse distributions, thus training a Bayesian neural network; at test time, this approach jointly inputs both the training and testing examples of the task without further optimization. The diversity of prior knowledge acquired in this manner has proven superior to standard tree-based techniques and confers notable advantages in the few-shot regime.

\subsection{Language-Interfaced Tabular Learning} An alternative direction for incorporating inductive biases and prior knowledge involves employing large language models (LLMs)~\cite{anil2023palm,openai2023gpt4}. As LLMs are pre-trained on massive, heterogeneous textual corpora, they exhibit a broad capacity for generalization and transfer to novel domains and tasks~\cite{brown2020language}. Current research has begun to exploit the rich background encoded by LLMs to enhance tabular learning. Typical methodologies involve translating tabular data structures into textual representations, which are then provided as part of LLM input prompts~\cite{dinh2022lift,hegselmann2023tabllm,wang2023anypredict}. Supplying feature and task descriptions alongside tabular entries within these prompts enables the inference of feature–task relationships for downstream prediction. Recent developments bifurcate into two predominant approaches: either adapting LLMs to tabular learning using parameter-efficient fine-tuning techniques—such as LoRA~\cite{hu2021lora} or IA3~\cite{liu2022few}—or harnessing in-context learning, whereby few-shot demonstrations are included within prompts and prediction proceeds without updating model parameters~\cite{dinh2022lift,hegselmann2023tabllm,nam2023semi,slack2023tablet}.

Although prior methods have significantly advanced few-shot tabular learning, their dependence on running inference through the LLM for every prediction introduces obstacles, particularly in industrial scenarios. This work seeks to move away from the typical black-box paradigm that processes each instance via an end-to-end LLM prediction. Rather, the proposed approach leverages the LLM to explicitly derive the rule-based conditions corresponding to each output class. \textsf{FeatLLM} translates these rules into executable code for feature engineering, generating new predictive features. This eliminates the necessity for continual inference with the LLM while harnessing in-context learning to induce class-specific rules informed by both prior knowledge and a limited number of examples.

\section{Methods}
\textsf{FeatLLM} is purpose-built to discover "rules"—class-specific conditions—based on the small set of examples provided, as illustrated in Figure~\ref{fig:main_model}. These rules are distilled through an LLM, then used to construct binary features for each data point, capturing whether each respective rule applies. The resulting feature set is passed through a dot product layer parameterized by non-negative, learnable weights to predict class probabilities. Training this system allows the model to identify which features most strongly influence the assignment of each class (see Section~\ref{Sec:training}). The entire method is iterated via bagging and ultimately aggregated through an ensemble strategy. The following describes the problem statement and offers detailed explanations for every stage of the process.

\vspace*{-0.12in}
\paragraph{Problem formulation.} Consider a tabular dataset composed of $N$ labeled examples, represented as $\mathcal{D} = \{(\mathbf{x}^i, \mathbf{y}^i)\}_{i=1}^{N}$. Each data point $\mathbf{x}^i$ is a $d$-dimensional feature vector, and the dataset $\mathcal{D}$ includes $d$ input variables with corresponding natural language names $F = \{f_j\}_{j=1}^{d}$, each paired with a provided definition. Under a classification framework, each label $\mathbf{y}^i$ belongs to a predefined class set. For $k$-shot learning scenarios, a subset of $k$ labeled examples (where $k < N$) is randomly drawn from $\mathcal{D}$ to train the model.

\subsection{Prompt Design for Extracting Rules}
\label{Sec:prompt_design}
To assist the LLM in extracting rules via more transparent and systematic reasoning, we encourage the approach to mirror the method an expert would take when interpreting tabular data. The constructed prompt consists of three essential parts (illustrated in Fig.~\ref{fig:prompt_for_rule} and detailed in Appendix~\ref{sec:example_rule}):

\vspace*{-0.12in}
\paragraph{Basic information description.} This component presents critical context necessary for the model to address the task at hand (highlighted as orange text in Fig.~\ref{fig:prompt_for_rule}). It covers a task description—including information about the labels—a presentation of the features, and a set of illustrative examples. The nature of the task is conveyed in the form of a question (for instance, ``Does this patient's myocardial infarction complications data indicate chronic heart failure? Yes or no?"). Additional examples are given in Appendix~\ref{sec:task_desc}. Feature descriptions specify each variable’s type and may optionally supply definitions; categorical variables may list possible categories. For demonstration, several training examples are converted into text using a structured serialization, displaying both the input and ground-truth label. This serialization method is consistent with prior work~\cite{dinh2022lift,hegselmann2023tabllm}:

\begin{align}
&\text{Serialize}(\mathbf{x}^i,\mathbf{y}^i, F) = \nonumber \\ 
&``\ f_1\ \text{is}\ \mathbf{x}^i_1.\ ... \ f_d\  \text{is}\  \mathbf{x}^i_d.\ \ \text{Answer:}\  \mathbf{y}^i\ ”
\end{align}

\vspace*{-0.12in}
\paragraph{Reasoning instruction.} Our objective is to improve the LLM's reasoning ability by supplying it with explicit reasoning instructions (refer to the blue highlights in Fig.~\ref{fig:prompt_for_rule}). The reasoning instruction commences with an opening line in the style of the chain-of-thought methodology~\cite{wei2022chain}. Subsequently, we segment the LLM's reasoning workflow into two distinct phases. First, the LLM is prompted to independently infer the underlying causal links or inclinations between the features and the task description, refraining from referencing sample demonstrations. Instead, it must draw upon general knowledge or common sense, enabling the comprehensive use of its pre-existing knowledge base about the problem context. During the second phase, the LLM examines the provided demonstrations alongside the knowledge gained in the initial step, synthesizing rules applicable to each class. This dual-phase procedure mitigates the risk of the model latching onto spurious associations from extraneous data columns and helps concentrate its focus on features with true significance. To control for situations where too few or too many rules are produced—leading to issues such as underfitting or unnecessarily large prompts—we enforce a cap on the number of rules extracted per class, set to 10\footnote{A detailed discussion about the optimal number of rules can be found in Section~\ref{sec:analysis}.}. Furthermore, as rules are articulated, the LLM is instructed to consult the feature type specified in the general information section, thus averting type inconsistency errors in rule formulation.

\vspace*{-0.12in}
\paragraph{Response instruction.} For ease of downstream parsing and utilization of generated rules, we steer the LLM’s output with targeted response instructions (see yellow highlights in Fig.~\ref{fig:prompt_for_rule}). The provided prompt spells out the expected template for answers within each class (i.e., the format of the response) as well as the standardized representation each rule should assume (i.e., the format of the rule). These specifications permit the LLM to choose appropriate structural formats contingent on the nature of each feature. Without further constraints, the LLM may connect several rules using logical connectors such as AND or OR. An illustrative example of the result is furnished in Appendix~\ref{sec:outcome-rule}.

\subsection{Inferring Class Likelihood via Rules}
\label{Sec:training}

\paragraph{Parsing rules for feature generation.} In this stage, the previously produced rules are employed to engineer new binary features. For every class, these features signal whether a given sample complies with the corresponding rules—yielding a value of '0' or '1'. Such features facilitate the estimation of a sample’s membership in each class. Nevertheless, as the rules formulated by the LLM are expressed in natural language, an effective parsing strategy must be devised for automated feature derivation. Even though rule formatting and response styles are regulated during the rule creation stage to streamline downstream parsing, the outputs from the LLM may still include unpredictable syntactic modifications~\cite{kovavc2023large}. For example, numeric expressions such as ``$>$” or ``$<$” can be represented verbosely as ``greater than” or ``less than”, complicating the parsing process.

To overcome these parsing obstacles, rather than developing intricate parsing software, we instead harness the LLM’s proficiency for interpreting semantic content amidst variable syntax. Moreover, since the LLM has been exposed to code during its training, it can generate straightforward programmatic solutions from textual guidelines. By incorporating the relevant function name, explanatory input and output descriptors, and the previously extracted rules within our prompt, we provide the LLM with sufficient context to perform the required transformation. Examples of such prompts and the resulting outputs are depicted in Figures~\ref{fig:prompt_for_parsing} and~\ref{sec:example_parsing}-\ref{sec:example_parse_outcome} in the Appendix. The code yielded by the LLM is subsequently executed using Python’s \texttt{exec()} routine to automate the data transformation process.

\vspace*{-0.12in}
\paragraph{Inferring class likelihood.} When presented with new binary features derived from class-specific rules, a straightforward approach to estimating the class likelihood for a given instance is to tally the number of satisfied rules for each class (i.e., sum the binary features by class). Nonetheless, the individual impact of rules and features can vary, making it essential to estimate their relative importance from the training data. To achieve this, we propose learning feature weights through a standard machine learning (ML) approach, applied to the binary feature vectors for each class. As an illustrative example, consider adopting a bias-free linear model. Define the binary feature vector extracted from the instance $\mathbf{x}^i$ for class $k$ as $\mathbf{z}_k^i \in \{0, 1\}^R$, where $R$ is the number of rules allocated per class and $c$ denotes the total number of classes. The resulting class probabilities $\mathbf{p}^i$ are determined by:

\begin{align}
\text{logit}_k^i &= \max(\mathbf{w}_k, 0) \cdot \mathbf{z}_k^i, \nonumber \\
\mathbf{p}^i &= \text{Softmax}({[\text{logit}_{1}^i, ..., \text{logit}_{c}^i]}). \label{eq:infer_prob}
\end{align}

Within this framework, $\mathbf{w}_k$ are the trainable parameters corresponding to feature weights in the linear model, capturing the relevance of each feature. By enforcing non-negative weights via projection onto the positive orthant, we prevent any feature generated by the LLM from negatively impacting the predicted class likelihood. These logits are then mapped to probability values for each class using the softmax operation. The parameters $\mathbf{w}_k$ are optimized using a cross-entropy loss function on a limited number of training examples.

\vspace*{-0.12in}
\paragraph{Ensembling with bagging.} As a final step, we repeat the overall method several times to obtain multiple inference models, and then aggregate their outputs to make the final class prediction via ensembling. Specifically, given a set of weight vectors trained over $T$ independent runs, indicated by $\{\mathbf{w}[t]\}_{t=1}^T$, we calculate the prediction for each sample by averaging the per-class probabilities $\mathbf{p}^i$—as defined in Eq.~\ref{eq:infer_prob}—across all models.

To inject randomness into rule generation within the ensemble framework, we configure the LLM inference temperature to a high value. Additionally, we randomize the sequence in which few-shot examples are presented, motivated by findings that such order can influence LLM output~\cite{lu2022fantastically}\footnote{An analysis regarding rule diversity is presented in Appendix~\ref{sec:diversity}}. To further leverage the variability in predictions for enhanced robustness, we employ bagging, in which subsets of features or instances are selected for each ensemble iteration; this approach proves especially beneficial when handling large numbers of training samples or feature sets. The ensemble strategy brings two key benefits. First, even if the LLM produces erroneous rules in certain runs, the correct rules derived in other runs may offset these, thereby bolstering framework robustness. This property aligns with the principle of LLM self-consistency~\cite{wang2022self}, since rules that recur across multiple ensemble trials are likelier to be reliable. Second, the ensemble method alleviates the limitation imposed by the LLM’s input prompt capacity. When tabular datasets surpass the prompt size limit, applying bagging reduces the input size, thus maintaining consistent model performance. While individual rule sets may not fully represent relationships involving every feature, the collective ensemble of many weak learners, each covering different subsets across trials, mitigates shortcomings from partial feature or instance representation.

\section{Experiments}
We assess the performance of \textsf{FeatLLM} on a range of tabular datasets and perform comprehensive analyses to better understand the behavior of our approach. Further results—including ablations over alternative LLMs, detailed qualitative evaluations, and additional studies—are included in the Appendix, owing to space limitations.

\subsection{Performance Evaluation}
\paragraph{Datasets.}
Our empirical evaluation covers 13 tabular datasets spanning binary and multi-class classification scenarios: (1) Adult~\cite{asuncion2007uci}, which predicts whether an individual’s annual income exceeds \$50,000; (2) Bank~\cite{moro2014data}, for assessing whether a bank customer will subscribe to a term deposit; (3) Blood~\cite{yeh2009knowledge}, aiming to identify donors likely to return; (4) Car~\cite{kadra2021well}, focused on evaluating car quality; (5) Communities~\cite{misc_communities_and_crime_183}, for estimating the crime rate of a specific locality; (6) Credit-g~\cite{kadra2021well}, used to classify individuals as good or bad credit risks; (7) Diabetes\footnote{\texttt{kaggle.com/datasets/uciml/pima-indians-diabetes-database}}, predicting diabetes diagnosis; (8) Heart\footnote{\texttt{kaggle.com/datasets/fedesoriano/heart-failure-prediction}}, targeting the identification of coronary artery disease; and (9) Myocardial~\cite{misc_myocardial_infarction_complications_579}, which determines the incidence of chronic heart failure in individuals.

We additionally introduce several recently released and synthetic datasets that have not been extensively examined in previous work and were absent from the pre-training corpus of large language models: (10) Cultivars, utilized for forecasting the grain yield of specific soybean cultivars\footnote{\texttt{archive.ics.uci.edu/dataset/913}}; (11) NHANES, aiming to classify the age group of individuals based on their records\footnote{\texttt{archive.ics.uci.edu/dataset/887}}; (12) Sequence-type, where the task is to determine whether a given sequence belongs to one of four categories: arithmetic, geometric, Fibonacci, or Collatz; and (13) Solution-mix, which requires predicting if the percentage concentration in a mixture derived from four component solutions exceeds 0.5. The initial two datasets were submitted to the UCI Machine Learning Repository after September 2023, thus ensuring their exclusion from the training data of the LLM API (gpt-3.5-turbo-0613) applied in this study. The latter pair consists of synthetic data, precluding the possibility of them being memorized by the LLM and thereby supporting robust evaluation.

The datasets differ considerably in terms of both scale and difficulty, as outlined in Table~\ref{Tab:basic-desc}.
Each provides comprehensive attribute names and definitions. In certain cases, such as the `Communities' and `Myocardial' datasets, the number of features surpasses one hundred, introducing difficulties due to the restricted context window of large language models.

\vspace*{-0.12in}
\paragraph{Baselines.}
We benchmark our methods against nine established baselines. The first set includes traditional approaches for tabular data: (1) Logistic regression (LogReg), (2) XGBoost~\cite{chen2016xgboost}, and (3) RandomForest~\cite{ho1995random}. The following two methods, (4) SCARF~\cite{bahri2022scarf} and (5) STUNT~\cite{nam2023stunt}, are designed for scenarios with access to unlabeled data, although acquiring large-scale unlabeled datasets can often be impractical in operational settings. Among the latest techniques is (6) TabPFN~\cite{hollmann2022tabpfn}, which leverages synthetic datasets sampled from broad distributions for model pretraining. The remaining three baselines harness the capabilities of LLMs: (7) In-context learning (In-context)~\cite{wei2022emergent}, (8) TABLET~\cite{slack2023tablet}, and (9) TabLLM~\cite{hegselmann2023tabllm}. In-context learning introduces a limited number of labeled examples directly into the model prompt without adjusting model weights. TABLET extends this approach by supplementing the prompt with auxiliary information such as rule-based and prototype-guided cues from an external classifier, aiming to enhance inference accuracy. Lastly, TabLLM incorporates parameter-efficient fine-tuning methods, such as IA3~\cite{liu2022few}, for adapting LLMs using small sample sets.

\vspace*{-0.12in}
\paragraph{Implementation details.}
\textsf{FeatLLM} employs GPT-3.5 as its LLM backbone, but the method itself does not depend on a specific LLM architecture (refer to Appendix~\ref{sec:backbones} for additional results using alternative backbones). The inference temperature for the LLM is configured at 0.5, while the top-p value retains the API default of 1. We fix the ensemble size to 20 and the number of extraction rules to 10. An examination of how hyper-parameters influence outcomes is provided in Figure~\ref{fig:hyperparameter2} and Appendix~\ref{sec:hyper-parameter}. For the linear model component, we optimize with Adam, applying a learning rate of 0.01, and train for 200 epochs. To select optimal epochs, $k$-fold cross-validation is applied. In terms of baseline reproduction, GPT-3.5 serves for in-context learning, while T0~\cite{sanh2022multitask} is used in the TabLLM experiments, aligning with original configurations. It is important to highlight that TabLLM restricts LLM usage to those with checkpoint access, which excludes GPT-3.5 from consideration. For conventional machine learning baselines such as LogReg, XGBoost, and RandomForest, their hyper-parameters are identified through a Grid-search procedure combined with $k$-fold cross-validation. The value of $k$ is set to either 2 or 4, ensuring the presence of at least one example per class in training splits. For complete implementation details, please consult Appendix~\ref{sec:baselines}.

\vspace*{-0.12in}
\paragraph{Results.}
Tables~\ref{tab:main1} and \ref{tab:main2} show comparative results over diverse datasets. Each experiment was conducted three times using varying random splits for training and testing samples, and we report the mean AUC and standard deviation across runs. Across these evaluations, \textsf{FeatLLM} repeatedly emerges as either the leading model or achieves the second highest rank relative to alternative baselines. Remarkably, our approach attains a level of performance with only 4 shots that is on par with traditional tabular models trained with 32 shots, as illustrated in Fig.~\ref{fig:compares}. It is also worth mentioning that although models like STUNT and TabPFN deliver noticeable improvements on certain datasets, their gains over established machine learning approaches are generally modest. Furthermore, methods employing LLMs, including in-context learning, TABLET, and TabLLM, encounter inference difficulties when applied to tabular data characterized by high dimensionality. In contrast, our method leverages bagging and ensembling strategies, which remain robust and effective for large-feature datasets and deliver consistent performance. It should also be acknowledged that the bagging and ensemble strategies developed here could be integrated with pre-existing LLM-based methods. However, doing so would entail doubling inference computations per sample, considerably raising computational demands and rendering such an approach inefficient.

\subsection{Analysis \& Discussion}
\label{sec:analysis}
\paragraph{Ablation study.} We conduct ablation experiments to investigate the individual impact of key framework components, specifically examining: (1) \textsf{-Tuning}: removing the weight tuning mechanism in the linear model and replacing it with a sum of binary features for logit calculation; (2) \textsf{-Ensemble}: excluding the ensembling stage; (3) \textsf{-Description}: not providing feature-level descriptions; and (4) \textsf{-Reasoning}: omitting Step-1 of the reasoning instructions during rule generation.

Table~\ref{tab:ablation} presents the mean AUC variation attributable to each ablation, aggregated over all datasets. Any ablation leads to a noticeable decrease in performance. Of particular note, the advantage conferred by weight tuning becomes more prominent as the shot count rises, suggesting that access to more data enhances the model’s capacity to accurately estimate the significance of each rule, thereby amplifying the value of the tuning process. Conversely, ablations involving feature descriptions and reasoning instructions have a greater negative effect in lower-shot scenarios, indicating that the ability to harness LLMs' prior knowledge is especially critical when labeled data are limited. Lastly, integrating the ensemble approach persistently yields performance gains across all experimental settings.

\vspace*{-0.12in}
\paragraph{Training \& inference time comparison.} We assess and contrast the training and inference times across various methods. All evaluations are conducted on the Adult dataset with a training set containing 16 examples. For all approaches, a single A100 GPU is utilized, with the exception of TabLLM, which is deployed across four A100 GPUs to enable model parallelism. Table~\ref{tab:cost} details the corresponding runtime metrics. Importantly, \textsf{FeatLLM} achieves a notably low inference latency, on par with traditional machine learning techniques. Conversely, inference with other LLM-centric strategies—including In-context learning, TABLET, and TabLLM—demands substantially greater computational time. Regarding training duration, \textsf{FeatLLM} performs similarly to other LLM-based frameworks, necessitating just 30 API calls. This minimal number of queries neither significantly impacts computational budget nor introduces substantial overhead, and may also be executed in parallel to improve efficiency.

\vspace*{-0.12in}
\paragraph{Quality of features generated by \textsf{FeatLLM}.}
To evaluate the utility of rule-based features produced by \textsf{FeatLLM} for practical applications, we conduct a straightforward analysis. Specifically, we calculate the average AUROC between each new feature and the target label, subsequently determining the proportion of features whose AUROC exceeds 0.5 for every dataset considered. The results are compiled in Table~\ref{tab:quality}. As indicated (see the ``Before tuning" column), most rules generated align closely with the target attribute and contribute meaningful information for solving the associated predictive task. Furthermore, throughout the process of calibrating a linear model with the available data, irrelevant or low-utility rules—defined as those with learned weights falling below a designated threshold—are systematically removed (refer to the ``After tuning" column). The threshold for this pruning is established at 0.1 based on an analysis of the distribution of learned weights.

\vspace*{-0.12in}
\paragraph{Effect of adding spurious correlations.} To evaluate the ability of different techniques to mitigate the influence of spurious correlations, especially under low-shot conditions, we carried out a focused analysis. This involved using the Heart dataset as our base, supplemented by the addition of randomly selected, unrelated columns from the Adult dataset, which have no bearing on the Heart dataset's prediction task. We systematically increased the number of these irrelevant columns and monitored the consequent performance variations of the tested models. To ensure comparability across methods, all models were provided with 8 labeled training samples, with no access to additional unlabeled data—thus precluding methods such as SCARF or STUNT from inclusion in this assessment.

In Figure~\ref{fig:spurious}, we display the changes in model performance as the prevalence of spurious correlations rises. Notably, \textsf{FeatLLM} exhibited the smallest decline in performance compared to competing baselines as more irrelevant features were introduced. This resilience is plausibly due to our framework's enforcement of prior knowledge to relate features to the specific task during rule extraction, which suppresses the derivation of rules involving uninformative columns and contributes to greater robustness. Our results show that rules constructed from these extraneous features constituted only about 1-2\% of the overall rule set. Nevertheless, it is important to mention that \textsf{FeatLLM} may still capture and act on spurious correlations—potentially conflating causal links between features and task—when additional columns originate from datasets with genuine relevance (see Appendix~\ref{sec:spurious-appendix} for further discussion).

\vspace*{-0.12in}
\paragraph{Hyper-parameter analysis}
\textsf{FeatLLM} relies on two principal hyper-parameters: the ensemble size and the number of rules to extract per class in each LLM interaction. We systematically evaluated the effect of these hyper-parameters on the model's efficacy. Our findings indicate that increasing the ensemble count consistently enhances performance; however, this improvement plateaus beyond approximately 20 ensembles (refer to Fig.~\ref{fig:appendix-hyperparameter1} in the Appendix). Regarding the number of rules per query, no statistically significant changes in performance were observed across different values, although setting it to 10 strikes an effective balance between prompt complexity and accuracy (see Fig.~\ref{fig:hyperparameter2}). We hypothesize that a larger set of rules augments the input dimensionality, injecting additional information, but at the risk of overfitting—especially apparent in the low-shot learning context.

\section{Conclusion}
This work introduces \textsf{FeatLLM}, which elevates few-shot tabular learning via LLMs to a new benchmark. \textsf{FeatLLM} distinguishes itself by using LLMs to articulate interpretable decision criteria akin to human feature engineering, rather than directly outputting predictions on a per-sample basis. By combining rule induction and bagged ensembling, \textsf{FeatLLM} operates with low inference costs, functions without gradient-based training—requiring only access via API—and sidesteps constraints related to feature dimensionality. As a result, \textsf{FeatLLM} demonstrates robust predictive power, presenting a highly competitive solution for employing LLMs in practical tabular tasks under data-scarce conditions.

At present, \textsf{FeatLLM} is specifically tailored to the low-shot learning environment. We envision future expansions where the framework will autonomously generate new features for larger datasets, incorporate a broader spectrum of feature construction mechanisms beyond rule-based approaches, and enhance the interpretability of its predictions, thereby broadening its applicability across a diverse range of domains.

\section{Broader Impact}
The \textsf{FeatLLM} framework leverages LLMs' foundational knowledge to extend the frontiers of tabular learning performance, surpassing conventional supervised methods. A core focus of our research is improving data efficiency, thus addressing overfitting risks endemic to limited labeled datasets by integrating external knowledge sources. \textsf{FeatLLM} is particularly valuable for application domains—such as finance and healthcare—where acquiring labeled data is both labor-intensive and expensive, and pretrained LLMs can contribute domain knowledge obtained from large relevant corpora. For \textsf{FeatLLM} to reach broader real-world deployment, it is crucial to study the effects of societal biases and misinformation encoded in LLMs’ prior knowledge, as these could disproportionately sway model outputs, potentially overshadowing empirical signals in scarce datasets.

\end{document}